%! AP Statistics — Unit 1, Lesson 1.10 — Homework KEY
\documentclass[10pt]{article}
\usepackage{apstats-article}
\usepackage{apstats-key}

%=============================================================================
\begin{document}

\pageheader{Unit 1, Lesson 1.10}{Homework — Key}

\vspace{0.10in}

\begin{tcolorbox}[colback=sky, colframe=navy, boxrule=0.8pt, arc=2mm,
    left=3mm, right=3mm, top=1.5mm, bottom=1.5mm]
\small Show all $z$-score calculations. Interpret every probability in context.
Use Table A or technology; indicate which you used.
\end{tcolorbox}

\vspace{0.12in}

% ════════════════════════════════════════════════════════════════════════════
% PART A — Empirical Rule
% ════════════════════════════════════════════════════════════════════════════
\boxguard
\begin{blurbbox}[Part A \quad The Empirical Rule]{navylight}
\small

IQ scores on a widely used test are approximately Normally distributed with
mean $\mu = 100$ and standard deviation $\sigma = 15$. Use the
68--95--99.7 Rule for all questions in this part (do not use Table A).

\vspace{0.10in}

\begin{enumerate}[leftmargin=1.5em, itemsep=8pt]

  \item Fill in the IQ values at each labeled point on the axis below.\\[6pt]
        \begin{center}
        \begin{tikzpicture}[scale=1.0]
          \draw[thick, navy]
            plot[domain=-3.5:3.5, samples=80] (\x, {2.2*exp(-\x*\x/2)});
          \draw[thick] (-3.8, 0) -- (3.8, 0);
          \foreach \z/\iq in {-3/55,-2/70,-1/85,0/100,1/115,2/130,3/145} {
            \draw (\z, -0.12) -- (\z, 0.12);
            \node[below, font=\small] at (\z, -0.15) {\ans{\iq}};
          }
          \node[below, font=\small\itshape] at (0, -0.55) {IQ score};
          \node[above, font=\small\bfseries\color{navy}] at (0, 2.22) {$\mu = 100$};
        \end{tikzpicture}
        \end{center}

  \item What percentage of people have IQ scores between 85 and 115?
        Between 70 and 130? Show reasoning.\\[5pt]
        Between 85 and 115: \ans{68}\% \quad Between 70 and 130: \ans{95}\%

  \item What percentage of people have IQ scores \textbf{above} 130?\\[5pt]
        \ansline{Above $\mu+2\sigma$: $(100\%-95\%)/2 = 2.5\%$ of people.}

  \item A person's IQ is reported as 145. Using the Empirical Rule, comment on how
        unusual this score is.\\[5pt]
        \ansline{$145 = \mu+3\sigma$. Since 99.7\% fall between 55 and 145, only about}\\[1pt]
        \ansline{0.15\% score above 145 --- an extremely unusual score.}

\end{enumerate}

\end{blurbbox}

\vspace{0.12in}

% ════════════════════════════════════════════════════════════════════════════
% PART B — z-scores and Normal table
% ════════════════════════════════════════════════════════════════════════════
\boxguard
\begin{blurbbox}[Part B \quad $z$-Scores and the Normal Table]{greenacc}
\small

Customer wait times at a coffee shop are approximately $N(4.5,\,1.2)$ minutes
(mean 4.5 min, SD 1.2 min).

\vspace{0.10in}

\begin{enumerate}[leftmargin=1.5em, itemsep=8pt]

  \item A customer waited 6 minutes. Calculate the $z$-score and interpret it.\\[5pt]
        $z = \ans{$(6-4.5)/1.2$} = \ans{$+1.25$}$\quad
        Interpretation: \ansline{6 min is 1.25 SD above the mean wait.}

  \item Use Table A to find the proportion of customers who wait \textbf{more than}
        6 minutes. Show all steps.\\[5pt]
        $P(\text{wait} > 6) = 1 - P(Z < \ans{1.25}) = 1 - \ans{0.8944} = \ans{0.1056}$

  \item Find the proportion of customers who wait \textbf{between 3 and 6 minutes}.\\[5pt]
        $z(3) = \ans{$(3-4.5)/1.2$} = \ans{$-1.25$}$\quad
        $z(6) = \ans{$(6-4.5)/1.2$} = \ans{$+1.25$}$\\[4pt]
        $P(3 < \text{wait} < 6) = P(Z < \ans{1.25}) - P(Z < \ans{-1.25}) = \ans{0.8944} - \ans{0.1056} = \ans{0.7888}$

  \item Find the wait time at the \textbf{10th percentile}.
        ($z^* \approx -1.28$ for the 10th percentile.)\\[5pt]
        $x = \mu + z^* \cdot \sigma = \ans{$4.5 + (-1.28)(1.2)$} = \ans{2.96}$ minutes

\end{enumerate}

\end{blurbbox}

\vspace{0.12in}

% ════════════════════════════════════════════════════════════════════════════
% PART C — AP-Style Free Response
% ════════════════════════════════════════════════════════════════════════════
\boxguard
\begin{blurbbox}[Part C \quad AP-Style Free Response]{redacc}
\small

Scores on a standardized graduate school entrance exam are approximately Normally
distributed with mean $\mu = 500$ and standard deviation $\sigma = 100$.

\vspace{0.08in}

\begin{enumerate}[leftmargin=1.5em, itemsep=8pt]

  \item What proportion of test-takers scored \textbf{above 650}?
        Show your $z$-score calculation, use Table A, and interpret the result.
\begin{work}
  z &= \dfrac{650-500}{100} = 1.50 \\
  P(Z<1.50) &= 0.9332 \\
  P(X>650)  &= 1-0.9332 = 0.0668
\end{work}
        \ansline{About 6.7\% of test-takers scored above 650.}

  \item What proportion scored \textbf{below 400}?
\begin{work}
  z &= \dfrac{400-500}{100} = -1.00 \\
  P(X<400) &= P(Z<-1.00) = 0.1587
\end{work}
        \ansline{About 15.9\% of test-takers scored below 400.}

  \item What score separates the \textbf{top 5\%} of test-takers from the rest?
        ($z^* \approx 1.645$ for the 95th percentile.) Show your work.
\begin{work}
  x &= \mu + z^*\sigma \\
    &= 500 + 1.645(100) \\
    &= 664.5 \approx 665
\end{work}
        \ansline{A score of about 665 separates the top 5\% from the rest.}

  \item A graduate program requires a score of at least 650 for admission.
        Based on your answer to (a), comment on how selective this requirement is.\\[5pt]
        \ansline{Very selective: only about 6.7\% of test-takers score at or above}\\[1pt]
        \ansline{650, so roughly 93\% of applicants fail this cutoff alone.}

\end{enumerate}

\end{blurbbox}

\end{document}
